\documentclass[a4paper,10pt]{report}\input{../0.Préambule/Préambule_cours_prof}\begin{document}

\pagestyle{empty}
\newpage
\hspace{8cm}\noindent NOM, Prénom et classe : \dotfill

\vspace{-1cm}
\section*{Nombres relatifs 2 \quad Sujet A}
\addcontentsline{toc}{section}{Nombres relatifs 2}

\noindent \textit{Il sera tenu compte dans la notation de la propreté ainsi que de la justification apportée à chacune des réponses. Cela représentera \textbf{1,5 points} sur l'ensemble du devoir}

\noindent \textit{L'usage de la calculatrice est \textbf{interdite}.} \hfill \textit{Durée : 55 minutes}

\paragraph{Exercice 1}: \textit{Exercice à compléter sur l'énoncé} \hfill \textbf{4.5 points}

Effectue les calculs suivants et complète sur l'énoncé
\begin{enumerate}
\begin{tabularx}{\linewidth}{*{3}{>{\arraybackslash}X}}
%\item $(-12) + (-15) = \dots ~\dots $ &
\item $(-20) + (+18) = \dots ~\dots $ &
\item $(-21) + (+21) = \dots ~\dots $ &
\item $(+24) + (-20) = \dots ~\dots $ \\
\item $(-2,1) + (+0,8) = \dots ~\dots $ &
\item $(-1,51) + (-0,14) = \dots ~\dots $ &
\item $(+0,3) + (-1) = \dots ~\dots $ \\
\item $(-1,15) + (+0,15) = \dots ~\dots $ &
\item $(+13,4) + (-14,9) = \dots ~\dots $ &
\item $(-12) + (+9,15) = \dots ~\dots $ \\
\end{tabularx}
\end{enumerate}

\paragraph{Exercice 2} : \textit{Exercice a faire sur une copie séparée .} \hfill \textbf{3 points}

\begin{enumerate}
\item A Spearfish, dans le Dakota du Sud, aux États-Unis la température était de $-20$\degre C et a gagné $27$\degre C en seulement $2$ minutes ! Quelle était la nouvelle température.
\item En France, à Mouth, le $13$ janvier $1968$, il y avait une température de $-36,7$\degre C et, en seulement quelques heures, la température a augmenté de $37,8$\degre C. Calcule cette nouvelle température.
\item Le $23$ mai $2015$, il faisait $-80,1$\degre C à Vostok en Antarctique alors qu'au Pakistan, le même jour il faisait $130$\degre C de plus ! Calcule la température au Pakistan le $23$ mai $2015$
\end{enumerate}

\paragraph{Exercice 3} : \textit{Exercice a faire sur une copie séparée .} \hfill \textbf{3 points}

Dans chaque cas, transforme la soustraction en addition puis effectue le calcul.
\begin{enumerate}
\begin{tabularx}{\linewidth}{*{3}{>{\arraybackslash}X}}
\item $(-21) - (+25)$ &
\item $(-52) - (-14)$ &
\item $(+42) - (+29)$ \\
\item $(-2,3) - (+2,4)$ &
\item $(-1,8) - (-2,5)$ &
\item $(-3,8) - (+5,8)$ \\
\end{tabularx}
\end{enumerate}

\paragraph{Exercice 4} : \textit{Exercice a faire sur une copie séparée .} \hfill \textbf{3 points}

Effectue les calculs suivants en regroupant les termes de même signes
\begin{enumerate}
\begin{tabularx}{\linewidth}{*{3}{>{\arraybackslash}X}}
\item $O = (-4) + (+6) + (-3)$ &
\item $C = (-15) + (-118) + (-47)$ \\
\item $E = (+1,8) + (-1,2) + (+3,4)$ &
\item $A = (-9) + (+13) + (+7) + (-11)$ \\
\item $N = (+1,9) + (+2,4) + (-8,6) + (+12,7)$ &
\item $S = (+8,92) + (+12) + (-8,92) + (-22)$ \\
\end{tabularx}
\end{enumerate}

\paragraph{Exercice 5}: \textit{Exercice a faire sur une copie séparée .} \hfill \textbf{3 points}

Pour chaque expression, effectue le calcul en regroupant les termes de même signe.
\begin{enumerate}
\begin{tabularx}{\linewidth}{*{2}{>{\arraybackslash}X}}
\item $A = -14 + 5 - 2 -2 - 23 + 33$ &
\item $S = 18 - 13 - 25 -0,8 + 2,7 - 3,7$ \\
\end{tabularx}
\end{enumerate}

\paragraph{Exercice 6}: \textit{Exercice a faire sur une copie séparée .} \hfill \textbf{3 points}

Dans un QCM de dix questions, une réponse juste rapporte $4$ points, une absence de réponse $0$ point et une mauvaise réponse enlève $3$ points.
\begin{enumerate}
\item Louis a $1$ bonnes réponses et $9$ mauvaises. Quelle est sa note ?
\item Quelle est la plus mauvaise note qu'il est possible d'obtenir à ce QCM ? la meilleur note ?
\item Nina a obtenu $14$ points. Donne une combinaison possible pour obtenir ce résultat.
\end{enumerate}

\paragraph{Exercice bonus}: \hfill \textbf{1 point}

Quatre boutons sont alignés comme le montre la figure : \quad\quad\bcsmbh \quad \bcsmmh \quad \bcsmbh \quad \bcsmmh

Deux d'entre eux ont une face joyeuse, et deux d'entre eux une face triste. Si on appuie sur une face, elle se transforme en son opposé (une face joyeuse devient triste et une triste devient joyeuse) et le ou les deux boutons de part et d'autres changent eux aussi d'expression.

\medskip
Quel nombre minimum de boutons faut-il presser pour obtenir uniquement des faces joyeuses ? Pour justifier votre réponse, vous pouvez représenter les faces au fur et à mesure qu'elles changent.


\newpage
\hspace{8cm}\noindent NOM, Prénom et classe : \dotfill

\vspace{-1cm}
\section*{Nombres relatifs 2 \quad Sujet B}
\addcontentsline{toc}{section}{Nombres relatifs 2 Sujet B}

\noindent \textit{Il sera tenu compte dans la notation de la propreté ainsi que de la justification apportée à chacune des réponses. Cela représentera \textbf{1,5 points} sur l'ensemble du devoir}

\noindent \textit{L'usage de la calculatrice est \textbf{interdite}.} \hfill \textit{Durée : 55 minutes}

\paragraph{Exercice 1}: \textit{Exercice à compléter sur l'énoncé} \hfill \textbf{4.5 points}

Effectue les calculs suivants et complète sur l'énoncé
\begin{enumerate}
\begin{tabularx}{\linewidth}{*{3}{>{\arraybackslash}X}}
%\item $(-12) + (-15) = \dots ~\dots $ &
\item $(-21) + (+21) = \dots ~\dots $ &
\item $(-20) + (+18) = \dots ~\dots $ &
\item $(+24) + (-20) = \dots ~\dots $ \\
\item $(+13,4) + (-14,9) = \dots ~\dots $ &
\item $(-1,15) + (+0,15) = \dots ~\dots $ &
\item $(-12) + (+9,15) = \dots ~\dots $ \\
\item $(-2,1) + (+0,8) = \dots ~\dots $ &
\item $(-1,51) + (-0,14) = \dots ~\dots $ &
\item $(+0,3) + (-1) = \dots ~\dots $ \\

\end{tabularx}
\end{enumerate}

\paragraph{Exercice 2} : \textit{Exercice a faire sur une copie séparée .} \hfill \textbf{3 points}

\begin{enumerate}
\item A Spearfish, dans le Dakota du Sud, aux États-Unis la température était de $-20$\degre C et a gagné $27$\degre C en seulement $2$ minutes ! Quelle était la nouvelle température.
\item Le $23$ mai $2015$, il faisait $-80,1$\degre C à Vostok en Antarctique alors qu'au Pakistan, le même jour il faisait $130$\degre C de plus ! Calcule la température au Pakistan le $23$ mai $2015$
\item En France, à Mouth, le $13$ janvier $1968$, il y avait une température de $-36,7$\degre C et, en seulement quelques heures, la température a augmenté de $37,8$\degre C. Calcule cette nouvelle température.
\end{enumerate}

\paragraph{Exercice 3} : \textit{Exercice a faire sur une copie séparée .} \hfill \textbf{3 points}

Dans chaque cas, transforme la soustraction en addition puis effectue le calcul.
\begin{enumerate}
\begin{tabularx}{\linewidth}{*{3}{>{\arraybackslash}X}}
\item $(-52) - (-14)$ &
\item $(-21) - (+25)$ &
\item $(+42) - (+29)$ \\
\item $(-2,3) - (+2,4)$ &
\item $(-3,8) - (+5,8)$ &
\item $(-1,8) - (-2,5)$  \\
\end{tabularx}
\end{enumerate}

\paragraph{Exercice 4} : \textit{Exercice a faire sur une copie séparée .} \hfill \textbf{3 points}

Effectue les calculs suivants en regroupant les termes de même signes
\begin{enumerate}
\begin{tabularx}{\linewidth}{*{3}{>{\arraybackslash}X}}

\item $P = (+1,8) + (-1,2) + (+3,4)$ &
\item $L = (-9) + (+13) + (+7) + (-11)$ \\
\item $A = (-4) + (+6) + (-3)$ &
\item $G = (-15) + (-118) + (-47)$ \\
\item $E = (+8,92) + (+12) + (-8,92) + (-22)$ &
\item $S = (+1,9) + (+2,4) + (-8,6) + (+12,7)$  \\
\end{tabularx}
\end{enumerate}

\paragraph{Exercice 5}: \textit{Exercice a faire sur une copie séparée .} \hfill \textbf{3 points}

Pour chaque expression, effectue le calcul en regroupant les termes de même signe.
\begin{enumerate}
\begin{tabularx}{\linewidth}{*{2}{>{\arraybackslash}X}}
\item $R = -14 + 5 - 2 -3 - 23 + 33$ &
\item $U = 18 - 13 - 24 -0,8 + 2,7 - 3,7$ \\
\end{tabularx}
\end{enumerate}

\paragraph{Exercice 6}: \textit{Exercice a faire sur une copie séparée .} \hfill \textbf{3 points}

Dans un QCM de dix questions, une réponse juste rapporte $4$ points, une absence de réponse $0$ point et une mauvaise réponse enlève $3$ points.
\begin{enumerate}
\item Louis a $3$ bonnes réponses et $7$ mauvaises. Quelle est sa note ?
\item Quelle est la plus mauvaise note qu'il est possible d'obtenir à ce QCM ? la meilleur note ?
\item Nina a obtenu $14$ points. Donne une combinaison possible pour obtenir ce résultat.
\end{enumerate}

\paragraph{Exercice bonus}: \hfill \textbf{1 point}

Quatre boutons sont alignés comme le montre la figure : \quad\quad\bcsmbh \quad \bcsmmh \quad \bcsmbh \quad \bcsmmh

Deux d'entre eux ont une face joyeuse, et deux d'entre eux une face triste. Si on appuie sur une face, elle se transforme en son opposé (une face joyeuse devient triste et une triste devient joyeuse) et le ou les deux boutons de part et d'autres changent eux aussi d'expression.

\medskip
Quel nombre minimum de boutons faut-il presser pour obtenir uniquement des faces joyeuses ? Pour justifier votre réponse, vous pouvez représenter les faces au fur et à mesure qu'elles changent.

\end{document}